\newcommand{\numalgorithms}{10}

\chapter{Algorithm description}

This chapter describes each of the \numalgorithms{} algorithms evaluated in this thesis,
grouped by method family. Table~\ref{tab:algo_overview}
lists them with their sensors and method families.

\begin{table}[H]
\centering
\caption{Evaluated algorithms overview}
\label{tab:algo_overview}
\begin{tabular}{clll}
\toprule
\# & Algorithm & Sensors & Method family \\
\midrule
1  & EKF Dead Reckoning    & Wheel + IMU        & Sensor fusion \\
2  & KISS-ICP              & LiDAR              & LiDAR odometry \\
3  & MOLA Odometry         & LiDAR              & LiDAR odometry \\
4  & Cartographer 2D       & LiDAR (2D scan)    & 2D LiDAR SLAM \\
5  & GMapping              & LiDAR (2D scan)    & 2D LiDAR SLAM \\
6  & SLAM Toolbox          & LiDAR (2D scan)    & 2D LiDAR SLAM \\
7  & RTABMap (LiDAR)       & LiDAR              & 3D LiDAR SLAM \\
8  & FAST-LIO2             & LiDAR + IMU        & LiDAR-IMU odometry \\
%9  & LIO-SAM               & LiDAR + IMU        & LiDAR-IMU SLAM \\
9  & cuVSLAM               & Stereo camera      & Visual SLAM \\
%10 & ORB-SLAM3 Stereo      & Stereo camera      & Visual SLAM \\
%11 & ORB-SLAM3 Stereo+IMU  & Stereo + IMU       & Visual-inertial SLAM \\
10 & RTABMap RGB-D         & Stereo camera      & Visual SLAM \\
\bottomrule
\end{tabular}
\end{table}

% ─────────────────────────────────────────────────────────────────────────────
\section{Sensor Fusion}

Sensor fusion combines wheel odometry and IMU to estimate pose without building a map.
It serves as a dead-reckoning baseline.

\subsection{EKF Dead Reckoning}

EKF (Extended Kalman Filter) node from the robot\_localization ROS2 package fuses wheel
odometry and IMU data to estimate the robot's pose without any geometric map.
Dead reckoning accumulates drift from wheel slip and noisy IMU data, so errors grow
unboundedly over time unless corrected.

% ─────────────────────────────────────────────────────────────────────────────
\section{LiDAR Odometry}

LiDAR odometry methods estimate the motion of the robot by aligning consecutive LiDAR scans
using the Iterative Closest Point (ICP) algorithm. Each new scan is compared to a local map of the previous few scans.
The best-fitting transformation is taken as the pose change.

These methods do not build a global map and do not detect loop closures, which
leads estimation error to accumulate proportionally to travelled distance.

LiDAR odometry is computationally efficient and works without visual features.
Methods that add loop closure or IMU fusion should outperform it 
especially on longer trajectories.

\subsection{KISS-ICP}

KISS-ICP~\cite{vizzo2023kissicp} (Keep It Small and Simple ICP) uses GICP with adaptive
parameters, requiring no manual tuning.
It was published in 2022 by researchers at the University of Bonn.

\subsection{MOLA Odometry}

MOLA~\cite{mola_docs} (Modular Optimization framework for Localization and mApping) was
developed at the University of Almería.
Only its LiDAR odometry component
(mola\_lidar\_odometry) was used here.

It uses point-to-plane ICP with an NDT voxel map.
Older scans are discarded as the robot moves to limit memory usage.
MOLA is built as a modular pipeline, which makes components easy to swap but requires more
configuration than KISS-ICP.

% ─────────────────────────────────────────────────────────────────────────────
\section{2D LiDAR SLAM}

2D LiDAR SLAM is designed for planar laser scanners that produce a single
horizontal sweep of range measurements. Representing the environment as a
horizontal plane reduces the pose estimation problem from 6-DoF to 3-DoF
($x$, $y$, $\theta$). This simplification lowers computational cost and is a
valid for wheeled robots on flat floors, but it discards all
height information. 2D SLAM is therefore unsuitable when full spatial reconstruction is needed or
when the robot operates on uneven terrain. The TurtleBot3 platform evaluated
in this thesis is a wheeled robot operating on a flat floor, so the planar
LiDAR is sufficient.

The simulation used in this thesis does not have a 2D laser scanner, only a point cloud from a
3D LiDAR. The lack of 2D LiDAR in simulation is caused by computational limitations.
Data from 3D LiDAR was processed offline: one horizontal ring was extracted
from each 3D scan and saved into a separate pre-processed bag. All three 2D
SLAM algorithms read from this bag.

Cartographer and SLAM Toolbox both build pose graphs and run global
optimisation when a loop closure is detected. GMapping uses a particle filter,
where each particle carries its own map.

\subsection{Google Cartographer 2D}

Cartographer 2D~\cite{hess2016cartographer} was developed by Google and uses a two-stage architecture.
The local stage matches each incoming scan against a small occupancy submap. 
Once a submap is complete it is frozen and
becomes a node in the pose graph. The global stage searches for loop closures
between frozen submaps using a branch-and-bound matcher, which is globally optimal.
Accepted loop closures are added as edges and the pose graph is then
optimised. The final map is a single occupancy grid assembled from all submaps.

\subsection{SLAM Toolbox}

SLAM Toolbox~\cite{macenski2021slamtoolbox} was developed by Steve Macenski and is the recommended 2D SLAM
solution for ROS\,2 Navigation2. Its defining feature is lifelong mapping:
the pose graph and map can be serialised to disk and resumed in a later
session, which makes it well suited to large environments that cannot be
mapped in a single run.

Scan matching uses a correlation-based matcher: a grid of candidate poses is
evaluated and the one that maximises the overlap score between the incoming
scan and the stored map is selected as the best alignment. Unlike Cartographer's
branch-and-bound approach, this can get stuck in local minima, making loop
closure more sensitive to tuning and environment geometry.

\subsection{GMapping}

GMapping~\cite{grisetti2007gmapping} is a Rao-Blackwellized particle filter SLAM.
There is no explicit loop closure. When the robot revisits a known area,
particle weights naturally shift toward poses consistent with the existing map,
which corrects the trajectory. This makes GMapping more tolerant of
repetitive environments where a correlation-based matcher might commit to a
false loop, but it also limits correction ability on long trajectories.

GMapping officially targets ROS\,1. The ROS\,2 port used here was compiled
from source against the openslam\_gmapping library.

% ─────────────────────────────────────────────────────────────────────────────
\section{3D LiDAR SLAM}

3D LiDAR SLAM builds maps from the full point cloud without discarding spatial
information, making it suitable when the environment cannot be described by a
single horizontal cross-section.

\subsection{RTABMap (LiDAR)}

RTABMap~\cite{labbe2019rtabmap} (Real-Time Appearance-Based Mapping) is a graph-based SLAM
system developed at Université de Sherbrooke. Pose estimates come from ICP-based scan
matching against a local point cloud map. The pose graph is globally optimised whenever
a loop closure is detected.

Loop closure detection is appearance-based. The system compares descriptors
(compact numerical summaries of the scan's geometric structure) of the current
scan against all stored nodes in the map, not just nearby ones. This means a
revisit can be detected even after a long detour, which methods that only search nearby nodes would miss.

To keep computation bounded, memory is split into Working Memory (WM) and Long-Term
Memory (LTM). When processing time exceeds a set threshold, the least-visited nodes
are moved to LTM. They are only retrieved when a loop closure candidate matches them.

% ─────────────────────────────────────────────────────────────────────────────
\section{Tightly-Coupled LiDAR-IMU}

These algorithms work with full 3D point clouds with an IMU fused directly
inside the same estimator.

The LiDAR provides geometric accuracy to estimate and remove the IMU's accumulated drift.
Neither sensor alone is sufficient.

% Two systems are evaluated. FAST-LIO2 is an odometry-only system with no loop closure.
% LIO-SAM adds loop closure on top of the same tight-coupling approach.
One system is evaluated: FAST-LIO2, an odometry-only system with no loop closure.

A rotating LiDAR produces points over a full revolution while the robot is
moving, so points captured at the start and end of the scan come from
slightly different poses.
Deskewing corrects this by transforming all points in a scan to a common
timestamp using IMU-interpolated poses.
IMU preintegration~\cite{forster2017manifold} avoids re-integrating all
high-frequency IMU measurements from scratch at every step by accumulating
them into a single compact motion estimate that remains valid during
optimisation.
FAST-LIO2 relies on both techniques.

\subsection{FAST-LIO2}

FAST-LIO2~\cite{xu2022fastlio2} was developed at the HKU-MARS laboratory. It
registers raw LiDAR points directly to the map. No feature extraction step is
needed. This makes it work well with LiDARs that have a small field of view or
non-repetitive scanning patterns, where feature-based methods fail.

It uses an iEKF: each point is matched to its nearest neighbours, a local plane is
fitted, and the point-to-plane distance drives the update, correcting pose and IMU
biases in one step.

The map is stored in an incremental k-d tree (ikd-Tree), which extends the
standard k-d tree (Section~\ref{sec:kdtree}) to support inserting and deleting
points without rebuilding the structure. This keeps search fast as the map
grows.

% \subsection{LIO-SAM}
%
% LIO-SAM~\cite{shan2020liosam} was developed at MIT and Stevens Institute of
% Technology. It uses a factor graph that combines IMU preintegration, LiDAR
% odometry, loop closure, and optionally GPS. When a loop closure is detected, the
% graph is globally optimised, correcting drift across the entire trajectory.
%
% Scan matching is feature-based. Each LiDAR scan is organised into rings,
% circular sweeps at different elevation angles. For each point, the algorithm
% checks how much the distances to neighbouring points on the same ring vary. On a
% flat surface the neighbouring distances are similar, so the point is classified
% as a planar feature. At an edge or corner the distances change sharply on one
% side, so the point is classified as an edge feature. These features are matched against a local map built from recent keyframes.

% ─────────────────────────────────────────────────────────────────────────────
\section{Visual SLAM}

Visual SLAM estimates the trajectory from image sequences using feature-based stereo vision.
The fundamental weakness is texture dependency: flat or uniformly coloured surfaces
produce few keypoints and the system loses tracking.

Three systems are evaluated, all using a stereo camera pair.

\subsection{cuVSLAM}

cuVSLAM~\cite{cuvslam_docs} is NVIDIA's visual SLAM library, distributed as part of the Isaac
ROS ecosystem. It runs the entire pipeline on an NVIDIA GPU, which allows dense feature
detection and tracking at rates that would be too slow on CPU. Stereo
matching produces depth for each tracked feature via triangulation. IMU measurements are
fused with the visual estimates to improve tracking stability and robustness to motion blur.
Pose increments are estimated from the tracked 3D feature positions and IMU data, and a
pose graph back-end with loop closure maintains global consistency. Because the library
is closed-source, internal parameters beyond those exposed through the ROS\,2 interface
are not accessible. The system requires a compatible NVIDIA GPU to run.

%\subsection{ORB-SLAM3 Stereo}
%
%ORB-SLAM3~\cite{campos2021orbslam3} is built on ORB-SLAM2 and used here in stereo
%configuration. The same ORB features are used for all tasks: tracking, mapping,
%relocalization, and loop closing. ORB are FAST corners extracted at multiple scale
%levels with a 256-bit descriptor. They are fast to compute and invariant to rotation
%and viewpoint changes.
%
%The system runs three parallel threads. The tracking thread computes the pose of each
%incoming frame by minimising reprojection error against the active map. The local
%mapping thread adds keyframes and culls redundant map points and keyframes. It refines
%the map by running local bundle adjustment over a local window of keyframes. The loop
%and map merging thread detects common regions between the active map and all stored
%maps using DBoW2. If the common region belongs to the active map, loop closing is
%performed and the Essential Graph is optimised. If it belongs to a different map, both
%maps are merged into a single active map.
%
%When tracking is lost, the system first tries to relocalize in all stored maps. If
%relocalization fails, the active map is stored and a new map is initialized from
%scratch.

\subsection{RTABMap RGB-D}

RTABMap~\cite{labbe2019rtabmap} in RGB-D mode reuses the same graph SLAM back-end as the LiDAR variant but
uses a visual front-end for loop closure. In this evaluation the depth image comes from stereo
matching on the Hawk camera. Wheel odometry provides the motion prior between frames. Loop
closure candidates are identified both by spatial proximity in the pose graph and by
bag-of-words matching of visual features extracted from the RGB image. Accepted loop
closures trigger a global graph optimisation that corrects the accumulated drift. The
accumulated map is a dense 3D point cloud built by projecting depth images into the world
frame at each pose.
% The same WM/LTM memory management applies.

%% ─────────────────────────────────────────────────────────────────────────────
%\section{Visual-Inertial SLAM}
%
%Visual-inertial SLAM extends stereo visual SLAM by adding an IMU to the estimator.
%The camera still provides depth through stereo triangulation, but the IMU contributes
%high-frequency rotation and acceleration measurements between frames. This addresses
%two weaknesses of stereo-only systems. Scale estimation is more reliable because the
%IMU constrains metric motion independently of the baseline. Brief tracking losses can
%be bridged by IMU prediction while the visual front-end searches for features to
%re-initialise.
%
%The IMU must run at a high rate for the preintegration to be accurate. Low IMU
%frequency causes the same problems as in the LiDAR-IMU case: inaccurate motion
%prediction and unreliable deskewing. The system also requires an initialisation phase
%where gravity direction, initial velocity, and IMU biases are estimated at the same
%time. Enough translational motion is needed to distinguish gravity from linear
%acceleration. A robot moving slowly in a straight line may fail to initialise.
%
%\subsection{ORB-SLAM3 Stereo+IMU}
%
%ORB-SLAM3 Stereo+IMU works like the stereo variant with IMU added to the estimator.
%The state is expanded to include body velocity and IMU biases. Inertial residuals are
%minimised jointly with visual reprojection errors using MAP estimation.
%
%Initialisation runs in three steps: vision-only, inertial-only, then joint
%visual-inertial MAP refinement. It achieves 5\% scale error in 2 seconds and 1\% in
%15 seconds. When tracking is briefly lost, the IMU predicts the body state and
%projects map points to search for matches and recover tracking.
